\scanpage{199}
\begin{center}- \mspage{58} -\end{center}
\margin{IX}

qui est une \uline{dualité parfaite}, i.e. identifiant chacun des deux facteurs du premier membre au dual de l'autre (au sens de 5.2.2). (NB les exposants ab désignent les ``parties abéliennes'', quotient de la partie fixe par la partie torique). Il est clair d'autre part que (5.3.2) est induit par l'accouplement
\[
(5.3.3)\qquad T_\ell(A^{\unsure{0|o}})^{\mathrm{ab}}\times T_\ell(A'^{\unsure{0|o}})^{\mathrm{ab}}\longrightarrow T_\ell(\mathbb G_{mS})
\]
déduit par passages aux quotients de l'accouplement (5.2.4.2). Celui-ci étant une dualité parfaite sur les fibres génériques, i.e. l'homomorphisme correspondant
\[
(5.3.4)\qquad T_\ell(A^{\unsure{0|o}})^{\mathrm{ab}}\longrightarrow D(T_\ell(A'^{\unsure{0|o}})^{\mathrm{ab}})
\]
étant un isomorphisme sur les fibres génériques, nous avons évidemment envie d'en déduire qu'il en est de même sur $S$ tout entier, ou ce qui revient au même, sur les fibres spéciales. Cette déduction est valable, trivialement, si $\ell\ne p$ ; elle est valable également, du moins dans le cas d'inégales caractéristiques, si $\ell=p$, grâce à un théorème profond de TATE [33, th.4]. Nous allons cependant pouvoir l'établir directement, sans l'aide du théorème de Tate, plus bas (7.4.6).

Pour préciser la signification géométrique du résultat que nous désirons établir, notons que l'accouplement déduit de (5.3.3) par passage aux fibres spéciales n'est autre que l'accouplement
\[
(5.3.5)\qquad T_\ell(B_{\unsure{o|0}})\times T_\ell(B'_{\unsure{o|0}})\longrightarrow T_\ell(\mathbb G_{mk})
\]
déduit de la biextension $W_{\unsure{o|0}}$ de $(A_{\unsure{o|0}}^{\unsure{0|o}},A'_{\unsure{o|0}}{}^{\unsure{0|o}})$ par passage au quotient, $B_{\unsure{o|0}}$ et $B'_{\unsure{o|0}}$
\margin{376}
\margin{376}
\margin{198}
